\subsection{Recommended compact companion: odd-parity calibration}
\label{sec:jones-compact}
The recommended compact companion has the same thirteen unrestricted integer
unknowns, exact degree $11\,543\,632$, and $518$ live binary operations
($326$ multiplications, $169$ additions, $23$ subtractions). Its literal DAG
is \texttt{generated/signed\_coprime\_compact\_split\_crt.json}.
The $520$-operation construction above is retained as the fully explicit
baseline. Only three source registers are changed:
\[
 \mathcal G_+=D_2(B)U_c^2+D_1(B)U_c+D_0(B)+B^{2161},\qquad
 L=B^{2279},\qquad \widehat N=\eta\epsilon B^{1201},
\]
with the same $\eta=2^{8\kappa-10}$. In the literal ratio registers use
\[
 U=\widehat N^2w^2,\qquad Y=\widehat N^2s.
\]
Everything else, including the mask, source targets, interval guard, CRT
registers, spectator stack, merger, and raw final norm, is unchanged.
The JSON register named \texttt{N} contains $\widehat N$, which can be
negative. \emph{Only in the proof}, instantiate the positive Jones parameter
by $N=|\widehat N|$. Since all uses of this register in the polynomial are
through its square, $N^2=\widehat N^2$ exactly. There is no absolute-value
gate or extra unknown in the equation.

\paragraph{Unconditional positivity.}
The same square completion, now with a slightly lower off-target monomial,
gives
\begin{align*}
 \mathcal G_+
 &>\tfrac14 B^{1929}U_c^2+B^{2161}
       -4A_*^2B^{2125}-2A_*B^{1719}\\
 &>\tfrac14 B^{1929}U_c^2+\tfrac12B^{2161},
\end{align*}
because $4A_*^2/B^{36}<1/4$ and $2A_*/B^{442}<1/4$. Thus
\[
 S_c>\tfrac12N_0B^{1929}U_c^2+U_c+N_0B^{2161}-\ell>0
\]
for every integer $g$, exactly as in the baseline proof. Hence $R>L^2$.
For the positive proof-side $N$ one has
\[
 N^2=2^{-21}a^{16}bB^{2402}
       <a^{15}B^{2403}<B^{2418}<B^{4558}=L^2.
\]
Also $N\ge c=bB^{116}$, since
$\eta B^{1085}\ge2|\epsilon|$, using $B\ge2|\epsilon|$ and $\eta\ge1$.
Finally $N\ge B^{1201}\ge8$. Thus all hypotheses $R\ge N\ge8$ and
$b,c\le N$ concern the positive number $|\widehat N|$, not signed
$\widehat N$. This is sufficient for every use of the ratio theorem and for
the canonical estimate proving $\Delta>0$.

\paragraph{A smaller block is sufficient.}
On the required range $0<g<c$, the previous bound
$2|\mathcal G|<B^{2162}/256$ still holds. Also
$T_*<B^{2162}/256$ and $2B^{2161}<B^{2162}/128$, since $B>256$.
Therefore
\[
 0<J_c<B^{2162}/64,\qquad
 S_c<B^{2279}/32,\qquad T_c<B^{2279}/64.
\]
Consequently $S_c+T_c<L=B^{2279}$. The pad is still above every target,
since $2161>2129$, and the same coefficient bound
$|h_e|<994A_*b^2<B/8$ covers its extra coefficient 1. The entire sparse
carry proof is therefore unchanged.

\paragraph{The exact square threshold uses the forced odd parity.}
After ratio soundness establishes $b=2^k$, the polynomial identity
$b=2\epsilon^2$ with integer $\epsilon\ne0$ forces $k$ odd and
$\epsilon^2=2^{k-1}$. Write $\nu=\kappa+2k$. For $L=B^{2279}$,
\begin{align*}
 \log_2L+s_2(T_c)+1
 &=2418\nu-31k-21\\
 &=2418\kappa+4805k-21\\
 &=2(8\kappa-10+1201\nu)+(k-1)
 =\log_2\widehat N^2.
\end{align*}
Thus $N^2=\widehat N^2$ has exactly the required central-binomial valuation
threshold. No even-$k$ assertion is needed: that branch cannot occur.
For completeness choose $\epsilon=2^r>0$ and $b=2^{1+2r}$, which are still
cofinal. The canonical $\omega$ remains a square. The same $b,c\le N$
estimate proves the interval guard. All later CRT, nonzero, square-recovery,
and sign-bootstrap arguments are identical. This proves the compact
companion's equivalence to a source zero and hence to a finite ZFC refutation.

\paragraph{Exact degree and accounting.}
Here the relevant degrees are
\[
 \deg a_0=261194,\quad \deg a_1=2634542,\quad
 \deg C_*=4874866,\quad\deg\Sigma=448475,
\]
and $\deg A_0=5771816$, $\deg A_1=3531492$. Again $A_0$ dominates
$X^2$ and $A_1$, and the leading term of the output is $A_0^2$.
Its diagonal coefficient is
\[
 630137\,a^{2885464}\eta^{960}\pmod{1000003}\ne0.
\]
This proves exact degree $11\,543\,632$, with $518$ operations under the
same gate metric. The change saves two operations and 2480 degrees against
the scaled split baseline, without changing the source tuple or number of
unknowns. The independent audit checks the all-integer source minimum,
compact block bounds, negative-$\epsilon$ normalization, odd-$k$ threshold,
and that the signed \texttt{N} node feeds only its square.
